% ############################################################################
% ##  SUPERSEDED 2026-08-03 -- NOT COMPILED BY ANY DOCUMENT. DO NOT EDIT.
% ##
% ##  This file is a per-theorem rendering from an ABANDONED architecture in
% ##  which paper/main.tex would \input one file per result.  paper/main.tex
% ##  DOES NOT EXIST.  The manuscript is paper/qmd/{main,supplement}.qmd, and
% ##  both \input the theory corpus directly.  Verified 2026-08-03: neither
% ##  .qmd references sections/theorems.
% ##
% ##  The text below is ROUND-2 vintage against a corpus now at dispatch 25.
% ##  66-78% of its substantive lines no longer match the corpus -- that is old
% ##  text the corpus has since revised, NOT unique work.  Every \label here is
% ##  also declared in the live corpus.
% ##
% ##  SOURCE OF RECORD:
% ##    statements  quality_reports/theory/panel-npiv-discrete/results.tex
% ##    assumptions quality_reports/theory/panel-npiv-discrete/assumptions.tex
% ##    proofs      quality_reports/theory/panel-npiv-discrete/proofs.tex
% ##
% ##  Kept rather than deleted because nothing in this repo is committed.
% ##  Full rationale: paper/sections/theorems/_SUPERSEDED.md
% ############################################################################

% ============================================================================
%  lem:adjoint -- closed form of the adjoint of the linearisation, with the
%  Bierens induction map that makes the representer evaluable at the data.
%  Project: panel-npiv-discrete.  Target of plan Sec.2.7; memo Sec.10 Step 3.
%
%  Source of record: quality_reports/theory/panel-npiv-discrete/{results,proofs}.tex
%  Assumptions cited by tag; see .../assumptions.tex.
%  Requires \IR, \un, \upmodels, \vdagger from preamble.tex.
%
%  DEPENDENCIES that must appear EARLIER in the compiled document:
%    def:model     -- eq:trialres, cited via prp:nophifree Step 0 in cor:rangegap;
%    def:spaces    -- V, L^2(Pi), nu_g (eq:nudef), the domain convention;
%    def:operator  -- T, M(u,t;varphi,theta) (eq:Mdef, eq:Tdef);
%    prp:wellposed -- Step 1's exhaustion argument, reused in lem:adjoint Step 2;
%    prp:frechet   -- A_varphi and eq:Adef;
%    prp:compact   -- the kernel k = nu_{g_{u,t}}/pi (eq:kernel), Steps 0-2;
%    lem:density   -- claim (i) and its Step 1, used for the density of R(W)
%                     and, with T replaced by T_1, inside lem:adjoint Step 4;
%    lem:bierenscore, lem:bierens -- Steps 2 and 5, used by prp:kerchar;
%    prp:nophifree -- Step 0 and eq:gproviso, used by cor:rangegap Step 1.
%  NOT a dependency, despite an earlier version of this header: def:admissible.
%  Conventions used: (P1), (P2), (P5), (P6a), (P7a), (P7b), and (P4c) in
%  prp:kerchar Step 5.  All are stated once in proofs.tex.
%
%  ASSUMPTION TAGS, BY RESULT rather than in aggregate (round-7 narrowing:
%  every "C1" below was found to consume only C1(ii); C1(i) has no consumer
%  before Step 6.  ROUND-8: four A1b(i) citations struck or redirected to the
%  DEFINITION of V -- see the round-8 note at the foot of this header):
%    lem:nusigma    C1(ii); A1b(ii) for (ii); A1 for (v).
%    lem:adjfubini  B1(i), (ii), (iv.c); C1(ii); C2'(i).  NOT A0: P is
%                   sigma-finite by def:model, and the round-6 appeal is struck.
%    lem:levels     A0 (X is a compact INTERVAL); C1(ii).
%    lem:adjoint    B1(i), (ii), (iv); C1(ii); C2'(i).   NOT B1(iii): the density
%                   half derives Leb(T_1) > 0 from B1(iv.b) instead.
%    cor:elldiff    B1(i), (ii), (iv.a); A2b (the PROPERTY tag, at a generic
%                   admissible representer).  A2'(i.a)-(i.b) is what specialises
%                   it to ell^Pi_star and is NOT its hypothesis.
%    cor:rangegap   the hypotheses of lem:adjoint; A1b(iv) via eq:gproviso.
%    prp:kerchar    (i) B1(i)-(iv), C1(ii), C2'(i);  (ii) adds A1 and A4(i).
%    cor:a1cprim    B1(i)-(iv), C1(ii), C2'(i), A1b(ii), A1d -- nothing more.
%
%  THIS FILE IS AN ALTERNATIVE RENDERING of EIGHT numbered environments --
%    lem:nusigma, lem:adjfubini, lem:levels, lem:adjoint, cor:elldiff,
%    cor:rangegap, prp:kerchar, cor:a1cprim -- plus the TWO remarks
%    rem:adjointreconcile and rem:ffcompare, TEN environments in all, carrying
%  the SAME \label keys as
%  quality_reports/theory/panel-npiv-discrete/results.tex.  The two must never
%  be \input into the same document.  (Round 7: the count word said NINE and the
%  list named eight; Part E carries thirteen environments in total, of which
%  rem-nophifree.tex renders the other three.  The partition was sound; only the
%  count was wrong.)
%
%  ROUND-8 SYNC.  THE HAND-KEPT EXCEPTION LIST IS ABOLISHED, and this is the
%  substantive change to this header.
%
%  The round-7 header enumerated the "known unsynced" items and warned, in its
%  own voice, that "a hand-kept list of exceptions is the SAME failure mode as a
%  hand-kept header, one remove further out -- it failed in exactly that way in
%  round 7, the list said FIVE and omitted a sixth".  IT THEN SAID SIX AND
%  OMITTED A SEVENTH: C25's rename of lem:nusigma(v)'s generic integrable
%  function from \vartheta to \Xi was live at four lines of this file, on the
%  letter that is BINDING on the Step-4 dispatch, and lem:solset(iii) closes by
%  \cref-ing this lemma "applied with \Xi = partial_u ell_star" -- so the
%  cross-reference already pointed across the gap.
%
%  A THIRD ITERATION OF THE SAME LIST IS NOT THE FIX.  The authoritative unsynced
%  set is whatever a MECHANICAL DIFF of this file against results.tex and
%  proofs.tex returns.  That diff is now a runnable artifact:
%
%      scripts/render-sync/check_rendering_sync.py
%
%  It compares, per \label key, the statement text of each environment rendered
%  here against the environment carrying the same key in
%  quality_reports/theory/panel-npiv-discrete/results.tex, and it sweeps a
%  declared symbol-rename table (glossary Sec.1.4-1.5) across the whole file.  It
%  reports; it does not edit.  RUN IT BEFORE ANY CLAIM THAT THIS FILE IS SYNCED.
%
%  WHAT ROUND 8 CLOSED IN THIS FILE, BY DOING IT RATHER THAN LISTING IT:
%    - \vartheta -> \Xi in lem:nusigma(v), statement and proof (C25);
%    - \mathcal B -> \Upsilon throughout cor:rangegap, statement and all three
%      proof steps (C20);
%    - \mathcal B(.) -> \mathfrak B(.) for the BOREL sigma-algebra at both sites,
%      lem:levels(ii) and the functional monotone class step of prp:kerchar (C20);
%    - A1b(i) -> the definition of V in lem:nusigma(ii), statement and proof, and
%      A1b(i) struck from lem:levels, cor:a1cprim and rem:ffcompare (auditor
%      Finding 1);
%    - cor:elldiff RESTATED at a generic admissible representer under the new
%      property tag A2b, with ell^Pi_star as the specialisation, statement and
%      proof, in the unstarred budgets Lambda^(0), Lambda^(1) (auditor Finding 4,
%      which gated Step 4);
%    - \mathrm{supp} -> \operatorname{supp} (spacing differs in the compiled
%      document).
%  WITH THE FIRST THREE DONE, THE \mathcal B COLLISION THAT WAS "LIVE IN THE
%  COMPILED DOCUMENT AND DEFERRED, NOT RESOLVED" IS RESOLVED.
%
%  WHAT REMAINS UNSYNCED IS NOT ENUMERATED HERE, ON PURPOSE.  Three round-7 proof
%  refinements (the (P6) integrable substitution lemma, the Pi_u (x) Pi_t Fubini
%  step, the Leb(U_k) attribution) and cor:elldiff's second continuity pass are
%  known to differ from proofs.tex; the diff script is what says whether that
%  list is complete, and this header is not.  Resync mechanically when Parts E
%  and F are generated together after Step 4.
%
%  SYNC DISCIPLINE.  This header was wrong in five ways after the round-5
%  resync, in a file whose mathematics was correct -- four environments missing
%  from the list above, two dependencies uncited, one cited but unused, two
%  conventions omitted, and the assumption tags given in aggregate rather than
%  by result.  A HAND-MAINTAINED HEADER IS THE SAME FAILURE MODE AS A
%  HAND-MAINTAINED RENDERING, one remove further out.  When Parts E and F are
%  generated together after Step 4, the label list, the dependency list and the
%  per-result tag list must all be produced MECHANICALLY from results.tex and
%  proofs.tex, and diffed, not retyped.
% ============================================================================

\begin{lemma}\label{lem:nusigma}
The period-difference operator is bounded with kernel the constants, and the two functionals it
generates differ in exactly one respect: the difference functional integrates to zero and the sum
functional does not.
Suppose Assumption C1(ii) holds and let $g$ be a random variable with $E|g\,h(X_1)|<\infty$ and
$E|g\,h(X_2)|<\infty$ for the $h$ named in each clause. Then
\begin{enumerate}
\item[(i)] $\|\Delta h\|_{L^{2}(P)}\le2\,\bar c_\pi^{1/2}\|h\|_{L^{2}(\pi)}$ for every $h\in\mathcal V$;
\item[(ii)] $\Delta h=0$ almost surely if and only if $h$ is almost surely constant, provided A1b(ii)
holds; consequently, under A1b(ii) alone, $\ker\Delta\cap\mathcal V=\{0\}$, the passage from
``constant'' to ``zero'' being $\int h\,d\pi=0$ for $h\in\mathcal V$, which is the definition of
$\mathcal V$ and not an assumption;
\item[(iii)] $E[g\,\Delta h]=\int_{\mathcal X}\nu_g(x)\,h(x)\,dx=\langle\nu_g/\pi,\,h\rangle_{L^{2}(\pi)}$,
for each $h$ satisfying the integrability proviso above --- in particular for every $h\in\mathcal V$ and
for $h\equiv1$;
\item[(iv)] $\int_{\mathcal X}\nu_g\,dx=0$ and $\int_{\mathcal X}\sigma_g\,dx=2\,E[g]$, where
\begin{equation}\label{eq:sigmadef}
\sigma_g(x)\;:=\;E[g\mid X_2=x]\,f_{X_2}(x)\;+\;E[g\mid X_1=x]\,f_{X_1}(x);
\end{equation}
\item[(v)] if in addition A1 holds and $g=\widetilde W\,\Xi(\Delta U,\Delta Z)$ for an integrable
$\Xi$, then $E[g]=0$ and hence $\int_{\mathcal X}\sigma_g\,dx=0$.
\end{enumerate}
\end{lemma}

\begin{proof}
Each clause is short; they are taken in order.

\emph{(i).} By $(a-b)^{2}\le2a^{2}+2b^{2}$ and C1(ii),
$E[(\Delta h)^{2}]\le2\int h^{2}(f_{X_1}+f_{X_2})\,dx\le4\bar c_\pi\|h\|^{2}_{L^{2}(\pi)}$ with
$\bar c_\pi=\max_j\|f_{X_j}/\pi\|_\infty$; take square roots.

\emph{(ii).} If $h$ is a.s.\ constant then $\Delta h=0$ a.s. Conversely $\Delta h=0$ a.s.\ says
$h(X_1)=h(X_2)$ a.s., so A1b(ii) makes $h$ a.s.\ constant; if $h\in\mathcal V$ then $\int h\,d\pi=0$ by
the definition of $\mathcal V$, and a constant with zero $\pi$-mean is zero because $\int\pi=1$.

\emph{(iii).} This is convention (P5), whose term-by-term integrability proviso is the hypothesis of the
lemma; the second equality divides and multiplies by $\pi$, which is bounded below away from zero on
$\mathcal X$ by C1(ii).

\emph{(iv).} Taking $h\equiv1$ in (iii) gives $\Delta h=0$, hence $\int\nu_g\,dx=0$. For $\sigma_g$,
$\int E[g\mid X_j=x]f_{X_j}(x)\,dx=E[g]$ for $j=1,2$ by the tower property, and the two terms add rather
than cancel. \emph{This one sign separates the first-order from the second-order analysis:} the
difference structure that makes $\int\nu_g=0$ automatic buys nothing about $\int\sigma_g$.

\emph{(v).} Under A1, $\Delta U\upmodels W\mid\Delta Z$, so
$E[\widetilde W\Xi(\Delta U,\Delta Z)\mid\Delta Z]
=E[\widetilde W\mid\Delta Z]E[\Xi\mid\Delta Z]=0$ by (P3); take expectations and apply (iv).
\end{proof}

\begin{lemma}\label{lem:adjfubini}
The $(u,t)$-integration defining the adjoint may be taken inside the expectation.
Suppose Assumptions B1(i), (ii), (iv.c), C1(ii) and C2$'$(i) hold. Then for every $\ell\in L^{2}(\Pi)$
and every $h\in\mathcal V$ the function $(u,t,\omega)\mapsto g_{u,t}(\omega)\Delta h(\omega)\ell(u,t)$ is
$\Pi\otimes P$-integrable, with
\begin{equation}\label{eq:adjenv}
\int_{\IR\times\mathcal T}E\big[\big|g_{u,t}\,\Delta h\big|\big]\,|\ell(u,t)|\,\Pi(du,dt)
\;\le\;2\,\bar c_\pi^{1/2}\,\|\bar f_0\|_\infty\,c_\Pi\,\|h\|_{L^{2}(\pi)}\,\|\ell\|_{L^{2}(\Pi)},
\end{equation}
and consequently $\int_{\IR\times\mathcal T}E[g_{u,t}\Delta h]\ell(u,t)\Pi(du,dt)
=E\big[\Delta h\int_{\IR\times\mathcal T}g_{u,t}\ell(u,t)\Pi(du,dt)\big]$.
\end{lemma}

\begin{proof}
The proof establishes the envelope \eqref{eq:adjenv}; Fubini's theorem then applies on $\Pi\otimes P$.

Joint measurability of $(u,t,\omega)\mapsto g_{u,t}(\omega)$ is Step~0 of the proof of
\cref{prp:compact}. By (P1), (P2) and $|f_{\Delta U\mid\cdots}(u)|\le\|\bar f_0\|_\infty$ from C2$'$(i),
$|g_{u,t}|\le\bar w(t)\|\bar f_0\|_\infty$, so by the Cauchy--Schwarz inequality in $P$ and
\cref{lem:nusigma}(i),
$E[|g_{u,t}\Delta h|]\le\bar w(t)\|\bar f_0\|_\infty(E[(\Delta h)^{2}])^{1/2}
\le2\bar c_\pi^{1/2}\|\bar f_0\|_\infty\bar w(t)\|h\|_{L^{2}(\pi)}$. Multiplying by $|\ell(u,t)|$ and
integrating against $\Pi$, the Cauchy--Schwarz inequality in $\Pi$ gives
$\int\bar w(t)|\ell|\,d\Pi\le(\int\bar w^{2}d\Pi)^{1/2}\|\ell\|_{L^{2}(\Pi)}=c_\Pi\|\ell\|_{L^{2}(\Pi)}$
by B1(iv.c). Combining the two displays gives \eqref{eq:adjenv}, which is finite; Fubini's theorem for
the $\sigma$-finite product $\Pi\otimes P$ permits the exchange.
\end{proof}

\begin{lemma}\label{lem:levels}
The orthogonality condition and the source condition are the same in the two parameterisations; the
remainder bound and the rate condition are not, and levels is the weaker of the two.
Suppose $\mathcal X$ is a compact interval as in A0, and Assumption C1(ii) holds; let
$\mathcal D^{-1}_0:L^{2}(\pi)\to\mathcal V$ be integration from the left endpoint of $\mathcal X$
followed by centring in $L^{2}(\pi)$. Then
\begin{enumerate}
\item[(i)] $\mathcal D^{-1}_0$ is bounded, with
$\|\mathcal D^{-1}_0\|\le2(|\mathcal X|\,\|1/\pi\|_\infty)^{1/2}$, and injective, while $\mathcal D$ is
unbounded on $L^{2}(\pi)$;
\item[(ii)] for integrable $\nu_g$, writing
$\mathcal N_g(v):=\int_{\{x\ge v\}\cap\mathcal X}\nu_g(x)\,dx$, one has $\mathcal N_g\equiv0$ if and only
if $\nu_g\equiv0$ almost everywhere on $\mathcal X$;
\item[(iii)] A2 holds for $\mathcal A_{\mathrm{cond}}$ if and only if the corresponding source condition
holds for $\mathcal A_{\mathrm{cond}}\mathcal D^{-1}_0$;
\item[(iv)] $\|h\|_{L^{2}(\pi)}\le\|\mathcal D^{-1}_0\|\,\|h'\|_{L^{2}(\pi)}$ whenever
$h=\mathcal D^{-1}_0h'$, so a rate in the derivative norm implies the same rate in the level norm, while
the converse fails.
\end{enumerate}
\end{lemma}

\begin{proof}
The four clauses are taken in order; the only substantive step is (ii), and (iii) follows from it by
injectivity rather than by a range computation.

\emph{(i).} Write $x_{\min}:=\min\mathcal X$, finite because $\mathcal X$ is compact, and
$\mathcal D^{-1}h'(x):=\int_{x_{\min}}^{x}h'(v)\,dv$. By the Cauchy--Schwarz inequality and C1(ii),
$|\mathcal D^{-1}h'(x)|^{2}\le|\mathcal X|\int_{\mathcal X}|h'|^{2}dv\le|\mathcal X|\|1/\pi\|_\infty\|h'\|^{2}_{L^{2}(\pi)}$
for every $x$, so $\|\mathcal D^{-1}h'\|^{2}_{L^{2}(\pi)}\le|\mathcal X|\|1/\pi\|_\infty\|h'\|^{2}_{L^{2}(\pi)}$
since $\int\pi=1$; centring adds at most the same amount again by the Cauchy--Schwarz inequality, giving
the stated constant. If $\mathcal D^{-1}_0h'=0$ then $\mathcal D^{-1}h'$ is $\pi$-a.e.\ constant, hence
constant by continuity, so $h'=0$ a.e.\ by Lebesgue differentiation. That $\mathcal D$ is unbounded is
exhibited by $h_k(x)=k^{-1}\sin(kx)$, for which $\|h_k\|_{L^{2}(\pi)}\to0$ while $\|h_k'\|_{L^{2}(\pi)}$
does not.

\emph{(ii).} One direction is immediate. Conversely, if $\mathcal N_g\equiv0$ then for $v_1<v_2$ the
difference gives $\int_{[v_1,v_2)\cap\mathcal X}\nu_g\,dx=0$; such intervals form a $\pi$-system
generating $\mathfrak B(\mathcal X)$ and containing $\mathcal X$, so the finite signed measure
$\nu_g\,dx$ vanishes identically and $\nu_g=0$ a.e.

\emph{(iii).} $(\mathcal D^{-1}_0)^{\ast}$ is injective on $\mathcal V$: if
$(\mathcal D^{-1}_0)^{\ast}\zeta=0$ then $\langle\mathcal D^{-1}_0h',\zeta\rangle=0$ for every $h'$, and
$\mathcal R(\mathcal D^{-1}_0)$ contains every centred $C^{1}$ function on $\mathcal X$, hence is dense
in $\mathcal V$, so $\zeta=0$. Since
$(\mathcal A_{\mathrm{cond}}\mathcal D^{-1}_0)^{\ast}=(\mathcal D^{-1}_0)^{\ast}\mathcal A^\ast_{\mathrm{cond}}$,
the derivative-side source condition reads
$(\mathcal D^{-1}_0)^{\ast}(\nu/\pi)=(\mathcal D^{-1}_0)^{\ast}\mathcal A^\ast_{\mathrm{cond}}\ell$, and
injectivity gives $\nu/\pi=\mathcal A^\ast_{\mathrm{cond}}\ell$, which is A2. The converse is immediate.

\emph{(iv).} The inequality is (i); the failure of the converse is the sequence $h_k$ of (i). \emph{The
consequence is that the rate requirement must be stated in the level norm, where it is the weaker of the
two demands.}
\end{proof}

\begin{lemma}\label{lem:adjoint}
The adjoint of the linearisation is the difference functional of the debiasing weight, divided by the
reference density, once the weight is read through the Bierens family.
Suppose Assumptions B1(i), (ii), (iv), C1 and C2$'$(i) hold. Then, with $g_{u,t}$ as in
\eqref{eq:kernel}:
\begin{enumerate}
\item[(i)] for every $\ell\in L^{2}(\Pi)$,
$(\mathcal A^\ast_{\mathrm{cond}}\ell)(x)=\int_{\IR\times\mathcal T}\mathsf k(u,t;x)\ell(u,t)\Pi(du,dt)\in\mathcal V$;
\item[(ii)] the induction map
\begin{equation}\label{eq:Wdef}
\big(\mathcal W\ell\big)(u,z)\;:=\;\varpi_u(u)\int_{\mathcal T}\ell(u,t)\,w\big(t'\Phi(z)\big)\,\Pi_t(dt)
\end{equation}
is a bounded linear map from $L^{2}(\Pi)$ into $L^{2}(P_{\Delta U,\Delta Z})$ with
$\|\mathcal W\|\le(\|\bar f_0\|_\infty\|\varpi_u\|_\infty\int_{\mathcal T}\bar w^{2}d\Pi_t)^{1/2}$, its
range is dense, and $\ker\mathcal W\subseteq\ker\mathcal A^\ast_{\mathrm{cond}}$;
\item[(iii)] for every $\ell\in L^{2}(\Pi)$,
\begin{equation}\label{eq:adjclosed}
\mathcal A^\ast_{\mathrm{cond}}\ell\;=\;\frac{\nu_{\widetilde W\,(\mathcal W\ell)}}{\pi},
\end{equation}
so $\mathcal A^\ast_{\mathrm{cond}}\ell=-\nu/\pi$ if and only if
$\nu_{\Delta Z+\widetilde W(\mathcal W\ell)}\equiv0$ almost everywhere on $\mathcal X$.
\end{enumerate}
\end{lemma}

\begin{proof}
The proof has four steps: the kernel form; boundedness of the induction map; the closed form, which is
one application of \cref{lem:adjfubini} followed by (P7a), the tower property and (P5); and the two
structural claims.

\emph{Step 1: part (i).}
By \cref{prp:compact}, $\mathcal A_{\mathrm{cond}}$ is the integral operator on $\mathcal V$ with kernel
$\mathsf k$ and $\iint|\mathsf k|^{2}\pi\,dx\,d\Pi<\infty$. Fubini, licensed by that
square-integrability and the Cauchy--Schwarz inequality, gives
$\langle\mathcal A_{\mathrm{cond}}h,\ell\rangle_{L^{2}(\Pi)}
=\langle h,\int\mathsf k(\cdot,\cdot;\cdot)\ell\,d\Pi\rangle_{L^{2}(\pi)}$, which identifies
$\mathcal A^\ast_{\mathrm{cond}}\ell$; membership of $\mathcal V$ is Step~2 of the proof of
\cref{prp:compact}.

\emph{Step 2: $\mathcal W$ is well defined and bounded.}
Joint measurability of $(u,z)\mapsto(\mathcal W\ell)(u,z)$ follows from joint measurability of
$(u,t,z)\mapsto\ell(u,t)w(t'\Phi(z))$ with the domination $|\ell(u,t)|\bar w(t)$ of (P2), by the
exhaustion over $\{\|t\|_1\le k\}$ of \cref{prp:wellposed} Step~1. By the Cauchy--Schwarz inequality in
$\Pi_t$ and (P2),
\begin{equation}\label{eq:Wpointwise}
\big|(\mathcal W\ell)(u,z)\big|^{2}\;\le\;\varpi_u(u)^{2}\Big(\int_{\mathcal T}\bar w^{2}d\Pi_t\Big)\int_{\mathcal T}\ell(u,t)^{2}\Pi_t(dt).
\end{equation}
Taking expectations over $(\Delta U,\Delta Z)$, using (P7b) and $\Pi_u(du)=\varpi_u(u)du$,
\[
E\big[(\mathcal W\ell)(\Delta U,\Delta Z)^{2}\big]
\le\Big(\int_{\mathcal T}\bar w^{2}d\Pi_t\Big)\int_{\IR\times\mathcal T}\ell(u,t)^{2}f_{\Delta U}(u)\varpi_u(u)\,\Pi(du,dt)
\le\|\bar f_0\|_\infty\|\varpi_u\|_\infty\Big(\int_{\mathcal T}\bar w^{2}d\Pi_t\Big)\|\ell\|^{2}_{L^{2}(\Pi)}.
\]

\emph{Step 3: the closed form, part (iii).}
By \cref{prp:compact} Step~1, $\mathcal A_{\mathrm{cond}}h(u,t)=E[g_{u,t}\Delta h]$. By
\cref{lem:adjfubini},
\begin{equation}\label{eq:adjstep}
\langle\mathcal A_{\mathrm{cond}}h,\ell\rangle_{L^{2}(\Pi)}
=E\Big[\widetilde W\,\Delta h\int_{\IR\times\mathcal T}w\big(t'\Phi(\Delta Z)\big)f_{\Delta U\mid\mathcal F_{XZW}}(u)\,\ell(u,t)\,\Pi(du,dt)\Big].
\end{equation}
Since $\Delta Z$ is $\mathcal F_{XZW}$-measurable, (P7a) with
$\zeta(u)=\int_{\mathcal T}\ell(u,t)w(t'\Phi(\Delta Z))\Pi_t(dt)$ --- $\Pi_u$-integrable by
\eqref{eq:Wpointwise} --- identifies the inner integral with
$E[(\mathcal W\ell)(\Delta U,\Delta Z)\mid\mathcal F_{XZW}]$ almost surely. \emph{This is the only step
that uses the absolute continuity of $\Pi_u$ in \textup{B1(iv.a)}, and without it the inner integral has
no representation as a function of $(\Delta U,\Delta Z)$.} Substituting into \eqref{eq:adjstep}, using
$\mathcal F_{XZW}$-measurability of $\widetilde W$ and $\Delta h$ and then the tower property,
$\langle\mathcal A_{\mathrm{cond}}h,\ell\rangle_{L^{2}(\Pi)}=E[\widetilde W(\mathcal W\ell)(\Delta U,\Delta Z)\Delta h]$.
Convention (P5) applies --- $E|\widetilde W(\mathcal W\ell)h(X_j)|\le\|\mathcal W\ell\|_{L^{2}}(E[h(X_j)^{2}])^{1/2}<\infty$
by (P1), Step~2 and C1(ii) --- and yields \eqref{eq:adjclosed}. The final claim follows because $\pi>0$
by C1(ii) and $g\mapsto\nu_g$ is linear.

\emph{Step 4: the structural claims of part (ii).}
If $\mathcal W\ell=0$ then $\widetilde W(\mathcal W\ell)=0$ a.s., so \eqref{eq:adjclosed} gives
$\mathcal A^\ast_{\mathrm{cond}}\ell=0$; hence
$\ker\mathcal W\subseteq\ker\mathcal A^\ast_{\mathrm{cond}}$. Density is shown in three parts.

\emph{A reduced generating family.} Put $\mathcal U_k:=\{u:\varpi_u(u)\ge1/k\}$,
$\mathcal T_1:=\mathcal T\cap\operatorname{supp}\Pi_t$, and let $\mathcal G_\ast$ be the collection of
$(u,z)\mapsto\varsigma(u)w(t_0'\Phi(z))$ with $\varsigma$ bounded measurable supported in some
$\mathcal U_k$ and $t_0\in\mathcal T_1$. If $\xi\perp\mathcal G_\ast$ then, since $\varpi_u>0$ everywhere
by B1(iv.a) gives $\bigcup_k\mathcal U_k=\IR$, dominated convergence with envelope
$|\xi|\|\varsigma\|_\infty\bar w(t_0)$ extends $E[\xi\varsigma(\Delta U)e^{t_0'\Phi(\Delta Z)}]=0$ to
every bounded measurable $\varsigma$, and Step~1 of the proof of \cref{lem:density} --- which requires of
the index set only positive Lebesgue measure, through \cref{lem:bierenscore} --- gives $\xi=0$.
\emph{The reduction is needed because B1(iv.a) makes $\varpi_u$ strictly positive pointwise but not
bounded below, so $\varpi_u^{-1}$ need not be integrable against an arbitrary $\varsigma\in C_b$.}

\emph{Index points.} By B1(iv.b) and the Lebesgue decomposition $\Pi_t$ has a nonzero absolutely
continuous part $\varpi_t\,d\mathrm{Leb}$ on $\mathcal T_0$ --- $\Pi_t$ itself is not assumed absolutely
continuous, and $\varpi_t$ is the density of its absolutely continuous part only. At every Lebesgue
density point $t_0$ of $\{\varpi_t>0\}$, $\Pi_t(B_\epsilon(t_0))>0$ for every $\epsilon>0$. The set of
such $t_0$ has positive Lebesgue measure and lies in $\mathcal T_1$.

\emph{The approximation.} Fix such a $t_0$ and $\varsigma$ supported in $\mathcal U_k$, and put
$\ell_\epsilon(u,t):=\varpi_u(u)^{-1}\varsigma(u)\Pi_t(B_\epsilon(t_0))^{-1}\un\{t\in B_\epsilon(t_0)\}$,
for which
$\|\ell_\epsilon\|^{2}_{L^{2}(\Pi)}=\Pi_t(B_\epsilon(t_0))^{-1}\int_{\mathcal U_k}\varsigma^{2}\varpi_u^{-1}du\le k\Pi_t(B_\epsilon(t_0))^{-1}\|\varsigma\|^{2}_{L^{2}}<\infty$.
Then $(\mathcal W\ell_\epsilon)(u,z)-\varsigma(u)w(t_0'\Phi(z))$ is $\varsigma(u)$ times a
$\Pi_t$-average of $w(t'\Phi(z))-w(t_0'\Phi(z))$ over $B_\epsilon(t_0)$, and for $w=\exp$ with
$\|\Phi(z)\|_\infty\le c_\Phi$ the mean value theorem bounds that difference by
$c_\Phi\|t-t_0\|_1e^{c_\Phi(\|t_0\|_1+\|t-t_0\|_1)}$ uniformly in $z$, hence by
$c_\Phi\epsilon\sqrt p\,e^{c_\Phi(\|t_0\|_1+\epsilon\sqrt p)}$ on $B_\epsilon(t_0)$. The bound tends to
zero with $\epsilon$, uniformly in $(u,z)$, so every element of $\mathcal G_\ast$ lies in
$\overline{\mathcal R(\mathcal W)}$ and the first paragraph gives density.
\end{proof}

\begin{corollary}[to \cref{lem:adjoint}]\label{cor:elldiff}
Regularity of a range-space representer transfers to its observable-space image, and the transfer
is what makes the sup-norm requirements of A2$'$ well posed.
Suppose Assumptions B1(i), (ii), (iv.a) hold and let $\ell^\Pi\in L^{2}(\Pi)$ be an
\emph{admissible representer} in the sense of A2b, with budgets $\Lambda^{(0)},\Lambda^{(1)}$. Then
$\ell:=\mathcal W\ell^\Pi$ is bounded, is continuously differentiable in $u$ with
\begin{equation}\label{eq:Wderiv}
\partial_u\ell(u,z)
=\varpi_u'(u)\!\int_{\mathcal T}\!\ell^\Pi(u,t)\,w(t'\Phi(z))\,\Pi_t(dt)
+\varpi_u(u)\!\int_{\mathcal T}\!\partial_u\ell^\Pi(u,t)\,w(t'\Phi(z))\,\Pi_t(dt),
\end{equation}
and satisfies $\|\ell\|_\infty\le\|\varpi_u\|_\infty\Lambda^{(0)}$ and
$\|\partial_u\ell\|_\infty\le\|\varpi_u'\|_\infty\Lambda^{(0)}+\|\varpi_u\|_\infty\Lambda^{(1)}$.
\emph{Specialisation.} Under A2$'$(i.a)--(i.b) the selected solution $\ell^\Pi_\star$ satisfies A2b with
the starred budgets, so the displays hold at $\ell=\ell_\star:=\mathcal W\ell^\Pi_\star$.
\emph{Nothing here uses that $\ell^\Pi_\star$ solves the adjoint equation}, which is why the statement is
available at a generic representer and why the results that apply it there may cite its own hypothesis
line rather than restate the bridge.
\end{corollary}

\begin{proof}
Fix $z$. By (P2), the integrand $u\mapsto\ell^\Pi(u,t)w(t'\Phi(z))$ is continuously differentiable
for $\Pi_t$-a.e.\ $t$ by A2b, with
\[
\sup_{u\in\IR}\big|\partial_u\ell^\Pi(u,t)\,w(t'\Phi(z))\big|\;\le\;
\Big(\sup_{u\in\IR}\big|\partial_u\ell^\Pi(u,t)\big|\Big)\,\bar w(t)
\qquad\text{for every }z,
\]
and the right-hand side is $\Pi_t$-integrable with integral $\Lambda^{(1)}<\infty$ \emph{by the
form of A2b, inherited from \textup{A2$'$(i.b)}, in which the supremum is inside the integral}. That
ordering is what makes
this a dominating function: a bound of the form $\sup_u\int(\cdot)\,d\Pi_t<\infty$ would give
$L^{1}(\Pi_t)$-boundedness in $u$ without uniform integrability, and Leibniz's rule requires the latter.
Differentiation under the integral sign therefore applies, uniformly in $z$; multiplying by $\varpi_u$,
continuously differentiable by B1(iv.a), and applying the product rule gives \eqref{eq:Wderiv}. The two
bounds follow from (P2) and the definitions of $\Lambda^{(0)}$, $\Lambda^{(1)}$.
\end{proof}

\begin{corollary}[to \cref{lem:adjoint}]\label{cor:rangegap}
The source condition and the orthogonality equation the estimator targets can differ only inside the
non-closedness of the range of the adjoint.
Suppose the hypotheses of \cref{lem:adjoint} hold and define the \emph{tilted difference map}
$\Upsilon:L^{2}(P_{\Delta U,\Delta Z})\to\mathcal V$ by $\Upsilon\Xi:=\nu_{\widetilde W\Xi}/\pi$.
Then $\Upsilon$ is a bounded linear operator with $\|\Upsilon\|\le2\bar c_\pi^{1/2}$,
$\mathcal A^\ast_{\mathrm{cond}}=\Upsilon\circ\mathcal W$, and
\begin{equation}\label{eq:rangegap}
\mathcal R\big(\mathcal A^\ast_{\mathrm{cond}}\big)\;\subseteq\;\mathcal R(\Upsilon)\;\subseteq\;
\overline{\mathcal R\big(\mathcal A^\ast_{\mathrm{cond}}\big)} .
\end{equation}
\end{corollary}

\begin{proof}
The proof checks that $\Upsilon$ is bounded, factors the adjoint through it, and reads the two
inclusions off density of $\mathcal R(\mathcal W)$ and continuity of $\Upsilon$.

\emph{Step 1: boundedness.} For $\Xi\in L^{2}(P_{\Delta U,\Delta Z})$ put
$g_\Xi:=\widetilde W\Xi(\Delta U,\Delta Z)$, so $E|g_\Xi|\le\|\Xi\|$ by (P1) and Cauchy--Schwarz and
$E|g_\Xi h(X_j)|\le\|\Xi\|(E[h(X_j)^{2}])^{1/2}<\infty$ for $h\in\mathcal V$ by C1(ii). By
\cref{lem:nusigma}(iii), (P1) and \cref{lem:nusigma}(i),
$|\int\nu_{g_\Xi}h\,dx|=|E[g_\Xi\Delta h]|\le\|\Xi\|\|\Delta h\|_{L^{2}(P)}\le2\bar c_\pi^{1/2}\|\Xi\|\|h\|_{L^{2}(\pi)}$.
\emph{Those two bounds are precisely clauses 1 and 2 of \eqref{eq:gproviso} for $g_\Xi$, so Step~0 of the
proof of \cref{prp:nophifree} applies verbatim and gives $\nu_{g_\Xi}/\pi\in L^{2}(\pi)$} --- which must
be established before the inner product below can be written, rather than read off it.
So $h\mapsto\int\nu_{g_\Xi}h\,dx$ is bounded on $\mathcal V$ with norm at most $2\bar c_\pi^{1/2}\|\Xi\|$;
since $\nu_{g_\Xi}/\pi\in L^{2}(\pi)$, the identity
$\int\nu_{g_\Xi}h\,dx=\langle\nu_{g_\Xi}/\pi,h\rangle_{L^{2}(\pi)}$ is well formed by C1(ii), so the
Riesz representer is the $\mathcal V$-projection of $\nu_{g_\Xi}/\pi$ --- and that projection is
$\nu_{g_\Xi}/\pi$ itself, since $\int(\nu_{g_\Xi}/\pi)\pi\,dx=\int\nu_{g_\Xi}\,dx=0$ by
\cref{lem:nusigma}(iv). Hence $\Upsilon\Xi=\nu_{g_\Xi}/\pi\in\mathcal V$.

\emph{Step 2: factorisation and the first inclusion.} \Cref{lem:adjoint}(iii) is
$\mathcal A^\ast_{\mathrm{cond}}\ell=\Upsilon(\mathcal W\ell)$, so
$\mathcal R(\mathcal A^\ast_{\mathrm{cond}})=\Upsilon(\mathcal R(\mathcal W))\subseteq\mathcal R(\Upsilon)$.

\emph{Step 3: the second inclusion.} Given $\Xi$, \cref{lem:adjoint}(ii) supplies $\ell_k$ with
$\mathcal W\ell_k\to\Xi$; continuity of $\Upsilon$ gives
$\mathcal A^\ast_{\mathrm{cond}}\ell_k\to\Upsilon\Xi$, so
$\Upsilon\Xi\in\overline{\mathcal R(\mathcal A^\ast_{\mathrm{cond}})}$. \emph{Only density of
$\mathcal R(\mathcal W)$ is used; the range is not characterised.}
\end{proof}

\begin{remark}\label{rem:adjointreconcile}
Four consequences of \cref{lem:adjoint} are worth separating from the statement.

\emph{(i) The representer is evaluable at the data, and that is now a theorem.} A solution of the adjoint
equation is an element of $L^{2}(\Pi)$, a function of $(u,t)$ with $t$ the Bierens index; A2$'$(i)
and \eqref{eq:JSig} evaluate
$\ell_\star(\Delta U,\Delta Z)$ with $z$ a value of $\Delta Z$. \emph{The citation read
``A2$'$(i)--(ii)'' through dispatch~13}; A2$'$(ii) is retired --- it is implied by A2$'$(i) in full
together with B1(i), (ii), (iv.a) --- and the clause that evaluates $\ell_\star$ at the data was always
(i), whose closing sentence fixes $\ell_\star:=\mathcal W\ell^\Pi_\star$ pointwise.
Part (ii) supplies the map between the
two and part (iii) shows it commutes with the adjoint equation, so the object A2$'$ and \eqref{eq:JSig}
speak about is constructed rather than posited. \emph{The second citation read ``A3'' through
dispatch~12}; that tag is retired and the display it presupposed is \eqref{eq:JSig}, which is what the
sentence always meant. The construction consumes exactly one thing beyond the round-2
assumption block, namely absolute continuity of the range measure's $u$-factor in B1(iv.a) --- a
restriction on a reference measure the analyst chooses, not on the data-generating process.

\emph{(ii) The correction measure is not free.} The map \eqref{eq:Wdef} integrates against $\Pi_t$, so
the correction measure of B1(v) is $\Pi_t$ and the first half of B1(v) is implied by B1(iv.c) through the
Cauchy--Schwarz inequality rather than assumed. This is consistent with the estimator solving for $L$
directly: by part (ii) the Bierens family and the sieve have the same closure, so no $t$-integral has to
be computed even though the population representer is defined through one.

\emph{(iii) What the construction does not give, and what it gives anyway.} $\mathcal W$ has dense range
but need not be surjective, so a selection criterion posed over all of $L^{2}(P_{\Delta U,\Delta Z})$ may
be unreachable; and $\mathcal W$ need not be injective, though
$\ker\mathcal W\subseteq\ker\mathcal A^\ast_{\mathrm{cond}}$ means no direction relevant to the selection
is lost. \emph{It does not follow that the gap between A2 and the equation the estimator targets is
uncontrolled.} By \cref{cor:rangegap} that gap is confined to
$\overline{\mathcal R(\mathcal A^\ast_{\mathrm{cond}})}\setminus\mathcal R(\mathcal A^\ast_{\mathrm{cond}})$,
\emph{with no characterisation of $\mathcal R(\mathcal W)$ required}. That set is nonempty under H1:
$\mathcal A^\ast_{\mathrm{cond}}$ is compact, being the adjoint of a compact operator, and a compact
operator has closed range if and only if that range is finite-dimensional, which H1's infinite-rank
clause excludes. \emph{Injectivity is neither available nor needed for this} --- the same paragraph
concedes $\ker\mathcal W\subseteq\ker\mathcal A^\ast_{\mathrm{cond}}$ two sentences above --- and an
earlier version gave injectivity as the reason. The gap is exactly where the ill-posedness lives.

\emph{(iii$'$) The derivative of the representer.} The sup-norm requirements on $\ell_\star$ remain
assumptions, but they are stated on $\ell^\Pi_\star$, in A2$'$(i.a)--(i.b), and \cref{cor:elldiff}
derives from them both \eqref{eq:Wderiv} and the two bounds. Rounds 1--2 assumed
$\|\partial_u\ell_\star\|_\infty<\infty$ for an object whose derivative nothing in the block made exist;
the $C^{1}$ clause of B1(iv.a) alone controls only the first term of \eqref{eq:Wderiv}. \eqref{eq:Wderiv}
is a \emph{conditional} identity whose hypothesis is A2$'$(i.b), a primitive awaiting L-21.

\emph{(iv) The $\nu/\pi$ referent is confirmed independently, and in the same form.} \emph{Notation:}
\citet{feve2014} write the endogenous regressors as $Z_1,Z_2$ with generic value $z$, the instrument as
$W$ with generic value $w$, and a generic element of the range space as $\psi$; the translation to this
paper is $Z_t\mapsto X_t$, $z\mapsto x$, and their $\psi$ is our $\ell$, the letter $\psi$ being reserved
here for the moment function. For the linear conditional-mean panel operator
$K\varphi=E[\varphi(Z_2)-\varphi(Z_1)\mid W]$, their Proposition~2.1 gives, \emph{verbatim in their
notation},
\[
\big(K^{\ast}\psi\big)(z)\;=\;\int\psi(w)\,\frac{f_{Z_2,W}(z,w)-f_{Z_1,W}(z,w)}{\pi(z)}\,dw ,
\]
followed immediately by the statement that $K^{\ast}\psi$ satisfies the normalising constraint
$\int[(K^{\ast}\psi)(z)]\pi(z)\,dz=0$. Translated, the numerator is $\nu_{\ell(W)}$ and the display is
$\nu_{\ell(W)}/\pi$. \emph{One difference of setting must be flagged rather than absorbed:} their $dw$
integral ranges over a general instrument space, whereas $W$ is binary here, so the counterpart in this
paper is the two-point sum
$\nu_{\ell(W)}(x)=\sum_{j\in\{0,1\}}\ell(j)\{f_{X_2,W}(x,j)-f_{X_1,W}(x,j)\}$; the
difference-of-densities structure is what carries over, not the integral. \emph{Both features of
\eqref{eq:adjclosed} therefore appear in print in the linear case:} the division by the reference
density --- with the same letter $\pi$ --- and the zero $\pi$-mean, which is membership of $\mathcal V$
and which \cref{prp:compact} Step~2 proves here from $\int\nu_g\,dx=0$ at $h\equiv1$. \emph{The
$\pi^{-1}$ is therefore not an artefact of the $L^{2}(\pi)$ convention adopted here; it is the shape the
adjoint takes in a first-differenced panel.} What transfers is that shape and nothing more: their $K$ is
a fixed linear conditional-expectation operator, while $\mathcal A_{\mathrm{cond}}$ is a Fr\'echet
derivative, and the induction map $\mathcal W$ has no counterpart there because their restriction is
already unconditional in the residual.
\end{remark}

\begin{proposition}\label{prp:kerchar}
A direction is annihilated by the linearised operator exactly when the instrument fails to move its
period difference, given the residual and the controls.
Suppose Assumptions B1(i)--(iv), C1 and C2$'$(i) hold, and let $h\in\mathcal V$. Then \emph{(i)} the
following are equivalent: \emph{(a)} $\mathcal A_{\mathrm{cond}}h=0$ in $L^{2}(\Pi)$; \emph{(b)}
$E[\widetilde W\{h(X_2)-h(X_1)\}\mid\Delta U,\Delta Z]=0$ almost surely. And \emph{(ii)} if in addition
Assumptions A1 and A4(i) hold, \emph{(b)} is equivalent to \emph{(c)}
$E[h(X_2)-h(X_1)\mid\Delta U,\Delta Z,W=1]=E[h(X_2)-h(X_1)\mid\Delta U,\Delta Z,W=0]$ almost surely.
\end{proposition}

\begin{proof}
Part (i) is Steps 1--4: (a)$\Rightarrow$(b) by stripping the Bierens index and then the $u$-index, and
(b)$\Rightarrow$(a) by reversing both. Part (ii) is Step 5, the extended binary identity (P4c). Steps
1--4 use neither A1 nor A4(i).

\emph{Step 1: (a) implies the pointwise family.}
The map $(u,t)\mapsto E[g_{u,t}\Delta h]$ is continuous in $u$ for each $t$ --- dominated convergence
with envelope $\bar w(t)\bar f_0|\Delta h|$, integrable by C2$'$(i), (P1), (P2) and
\cref{lem:nusigma}(i), together with continuity of $u\mapsto f_{\Delta U\mid\mathcal F_{XZW}}(u)$ from
C2$'$(i) --- and real-analytic in $t$ by Step~2 of the proof of \cref{lem:bierenscore}. Step~5 of the
proof of \cref{lem:bierens}, which consumes only B1(iv.a) and B1(iv.b), upgrades the $\Pi$-a.e.\
statement to $\mathcal A_{\mathrm{cond}}h(u,t)=0$ for every $(u,t)\in\IR\times\mathcal T$.

\emph{Step 2: stripping the Bierens index.}
Fix $u$, set $\varepsilon_u:=\widetilde Wf_{\Delta U\mid\mathcal F_{XZW}}(u)\Delta h$, integrable since
$E|\varepsilon_u|\le\|\bar f_0\|_\infty E|\Delta h|<\infty$, and note
$\mathcal A_{\mathrm{cond}}h(u,t)=E[\varepsilon_ue^{t'\Phi(\Delta Z)}]$. By \cref{lem:bierenscore} with
$V=\Delta Z$, under B1(i)--(iii), $E[\varepsilon_u\mid\Delta Z]=0$ almost surely for each $u$.

\emph{Step 3: stripping the $u$-index.}
For $\varsigma\in C_b(\IR)$, multiply by $\varsigma(u)$ and integrate over $u\in\IR$. The interchange is
licensed by Fubini with the envelope
$|\varsigma(u)\widetilde Wf_{\Delta U\mid\mathcal F_{XZW}}(u)\Delta h|\le\|\varsigma\|_\infty f_{\Delta U\mid\mathcal F_{XZW}}(u)|\Delta h|$,
whose double integral is $\|\varsigma\|_\infty E|\Delta h|<\infty$ by (P1),
$\int_\IR f_{\Delta U\mid\mathcal F_{XZW}}(u)du=1$ and \cref{lem:nusigma}(i). Then (P6a)'s substitution
lemma together with the tower property, $\widetilde W\Delta h$ being $\mathcal F_{XZW}$-measurable and
$\sigma(\Delta Z)\subseteq\mathcal F_{XZW}$, gives
\begin{equation}\label{eq:kerfam}
E\big[\widetilde W\,\Delta h\,\varsigma(\Delta U)\,\big|\,\Delta Z\big]=0\quad\text{a.s., for every }\varsigma\in C_b(\IR).
\end{equation}
Multiplying by a bounded $\beta(\Delta Z)$ and taking expectations, the functional monotone class
theorem --- the product functions $\varsigma\otimes\beta$ are closed under multiplication, contain the
constants and generate $\mathfrak B(\IR)\otimes\mathfrak B(\IR^{p})$, and the dominating function is
$|\widetilde W\Delta h|\sup_k\|\Lambda_k\|_\infty$ --- gives $E[\widetilde W\Delta h\,\Lambda]=0$ for
every bounded $\sigma(\Delta U,\Delta Z)$-measurable $\Lambda$, which is (b).

\emph{Step 4: (b) implies (a).}
(b) recovers \eqref{eq:kerfam}, which says the finite signed measure
$A\mapsto E[\widetilde W\Delta h\un\{\Delta U\in A\}\mid\Delta Z]$ integrates every $\varsigma\in C_b$ to
zero, hence vanishes, hence has zero density $E[\varepsilon_u\mid\Delta Z]=0$ for a.e.\ $u$ and, by the
continuity of Step~1, for every $u$. Multiplying by $e^{t'\Phi(\Delta Z)}$ and taking expectations gives
$\mathcal A_{\mathrm{cond}}h(u,t)=0$ for every $(u,t)$.

\emph{Step 5: part (ii) --- (b) is equivalent to (c), and here A1 and A4(i) enter.}
By A1, $P(W=1\mid\Delta U,\Delta Z)=\varrho(\Delta Z)$, so (P4c) with
$\mathcal C=\sigma(\Delta U,\Delta Z)$ and $\Lambda=\Delta h$ gives
$E[\widetilde W\Delta h\mid\Delta U,\Delta Z]=\varrho(1-\varrho)\{E[\Delta h\mid\Delta U,\Delta Z,W=1]-E[\Delta h\mid\Delta U,\Delta Z,W=0]\}$,
both cell-conditional expectations being well defined because A4(i) makes both $W$-cells non-null. A4(i)
also gives $\varrho(1-\varrho)\ge\varepsilon_\varrho(1-\varepsilon_\varrho)>0$, so the product vanishes
if and only if the brace does.
\end{proof}

\begin{corollary}[to \cref{prp:kerchar}(i)]\label{cor:a1cprim}
Instrument completeness given the residual, together with measure separability, delivers first-order
local identification.
Suppose Assumptions B1(i)--(iv), C1, C2$'$(i), A1b(ii) and A1d hold. Then A1c holds, that is,
$\ker\mathcal A_{\mathrm{cond}}=\{0\}$.
\end{corollary}

\begin{proof}
Let $h\in\ker\mathcal A_{\mathrm{cond}}\subseteq\mathcal V$. By \cref{prp:kerchar}(i), clause (b) holds:
$E[\widetilde W\Delta h\mid\Delta U,\Delta Z]=0$ almost surely; and $\Delta h$ lies in
$L^{2}(P_{X_1,X_2})$ by \cref{lem:nusigma}(i). A1d gives $\Delta h=0$ almost surely, and
\cref{lem:nusigma}(ii), which consumes A1b(ii) and the definition of $\mathcal V$, gives $h=0$.
\emph{The corollary reads clause
\textup{(b)}, not clause \textup{(c)}, so it uses neither A1 nor A4(i):} A1d is a statement about the
operator $\Lambda\mapsto E[\widetilde W\Lambda\mid\Delta U,\Delta Z]$, and so is clause (b).
\end{proof}

\begin{remark}\label{rem:ffcompare}
\Cref{cor:a1cprim} is the transplant of \citet[Prop.~2.2]{feve2014} to the conditional-independence
design, and it is worth saying what the transplant delivers and what it does not.

\emph{Notation, fixed once.} \Citet{feve2014} write the endogenous regressors as $Z_1,Z_2$ and the
instrument as $W$; here the endogenous regressors are $X_1,X_2$ and $Z$ is reserved for the exogenous
controls. All quotations below are verbatim \emph{in their notation}; the translation is
$Z_t\mapsto X_t$, and no symbol inside a quotation has been altered.

\emph{What it delivers.} Their Proposition~2.2 states: ``Let $\varphi\in E$. If $(Z_1,Z_2)$ is strongly
identified by $W$ and if $Z_1$ and $Z_2$ are measurably separated the function $\varphi$ is identified.''
The two primitives become A1d and A1b(ii) here, and the argument is the same in outline: completeness
removes the period difference from inside the conditional expectation, and measurable separability then
forces the direction to be constant, which the definition of $\mathcal V$ sets to zero because the
direction lies in $\mathcal V$. This supplies A1c with a primitive
sufficient condition it did not previously have.

\emph{What it does not deliver.} It does not reach global identification of $\vdagger$. In
\citet{feve2014} the restriction is a conditional mean and the operator is linear in $\varphi$, so
injectivity of the operator \emph{is} identification. Under conditional independence the operator
equation is nonlinear in $\varphi$ and the two statements separate: membership of $\mathcal V_{CI}$ is the
family $E[\widetilde W\varsigma(\Delta Y-\Delta\varphi)\mid\Delta Z]=0$ indexed by test functions of the
\emph{observed} residual at the trial value, whereas running the argument of \cref{prp:kerchar} globally
would need the strictly larger family
$E[\widetilde W\varsigma(\Delta Y-\Delta\varphi)\varkappa(\Delta U)\mid\Delta Z]=0$, indexed also by test
functions of the \emph{true} residual, about which the restriction says nothing at a trial $\varphi$. The
two coincide at $\varphi=\vdagger$ and nowhere else. \Cref{prp:compact} closes off the obvious repair:
because $\mathcal A_{\mathrm{cond}}$ is compact it is not bounded below, so A1c does not deliver even a
neighbourhood of uniqueness.

\emph{Why the gap is recorded rather than closed, and the precedent is sharper than a difficulty claim.}
\citet{feve2014} reject the hypothesis of their own Proposition~2.2, and their stated reason is the
excess A1d reproduces exactly: ``The assumption of Proposition 2.2 is too strong because it assumes that
$E(g(Z_1,Z_2)\mid W)=0$ implies $g(Z_1,Z_2)=0$ a.s.\ for any $g$, and we only use this property for
$g(Z_1,Z_2)$ of the form $\varphi(Z_2)-\varphi(Z_1)$.'' A1d has the same shape: it quantifies over every
$\Lambda\in L^{2}(P_{X_1,X_2})$, while \cref{cor:a1cprim} consumes it only at $\Lambda=\Delta h$, a
difference. \emph{Both papers reach a completeness condition stated over a class strictly larger than the
one the argument uses, and both decline to repair it by restricting the class}, because injectivity of
$\Lambda\mapsto E[\widetilde W\Lambda\mid\Delta U,\Delta Z]$ on differences alone is no longer a
recognisable completeness condition and has no primitive characterisation of its own.

\emph{The two asymmetries are not the same asymmetry.} Theirs runs from the assumption to its use: the
hypothesis is \emph{stronger than the argument consumes}. The one in the previous paragraph runs the
other way: the restriction supplies \emph{less information than the argument needs}, the $\varsigma$
family against the $(\varsigma,\varkappa)$ family. They are opposite inequalities, and they meet at one
point --- \emph{the recognisable completeness conditions available in either design are stated over a
coarser class than the design's own information}, so an assumption calibrated to that class is too
strong and the information calibrated to it is too weak. \emph{That is where a future improvement would
have to act, in either paper}: not by strengthening the completeness condition, but by finding a
restriction on $P$ stated over exactly the class the argument uses. A1b(iii) is accordingly kept as an
assumption, with A1d recorded as the route that reaches the linearised condition and stops there, and
with its own excess named rather than hidden.
\end{remark}
